%! Author = Omar Iskandarani
%! Date = 3/7/2026
%! Affiliation = Independent Researcher, Groningen, The Netherlands
%! License = © 2026 Omar Iskandarani. All rights reserved. This manuscript is made available for academic reading and citation only. No republication, redistribution, or derivative works are permitted without explicit written permission from the author. Contact: info@omariskandarani.com
%! ORCID = 0009‑0006‑1686‑3961
%! DOI = 10.5281/zenodo.yyy

\newcommand{\paperdoi}{10.5281/zenodo.yyy}
\newcommand{\papertitle}{\textbf{Hydrodynamics of Swirl‑Strings:}\ \Large Emergent Quantum Constants, Topological Leptons and Discrete Scale Invariance}
\newcommand{\paperkeywords}{Swirl-String Theory, hydrodynamics, vortex knots, electron mass, Planck constant, Bohr magneton, lepton generations, discrete scale invariance}

%=========================================
% PREAMBLE, PACKAGES AND DOCUMENT CONFIGURATION
%=========================================
\documentclass[11pt]{article}

\usepackage{amsmath,amssymb,amsfonts,bm}
\usepackage{siunitx}
\usepackage[hidelinks]{hyperref}
\usepackage[a4paper,margin=1in]{geometry}
\usepackage[T1]{fontenc}
\usepackage[utf8]{inputenc}
\usepackage{booktabs}
\usepackage{titlesec}
\usepackage[absolute,overlay]{textpos}

% swirl arrows (context‑aware) and canonical macros
\newcommand{\swirlarrow}{\mkern-2mu\scriptscriptstyle\boldsymbol{\circlearrowleft}}
\newcommand{\vswirl}{\mathbf{v}_{\mkern-2mu\scriptscriptstyle\boldsymbol{\circlearrowleft}}}
\newcommand{\SwirlClock}{S_{(t)}^{\mkern-2mu\scriptscriptstyle\boldsymbol{\circlearrowleft}}}
\newcommand{\Fmaxswirl}{F^{\max}_{\mkern-1mu\scriptscriptstyle\boldsymbol{\circlearrowleft}}}
\newcommand{\Fmax}{F^{\max}_{\mkern-1mu\scriptscriptstyle\boldsymbol{\circlearrowleft}}}
\newcommand{\FmaxEM}{F^{\max}_{\mathrm{EM}}}
\newcommand{\FmaxG}{F_{\mathrm{G}}^{\max}}               % G‑like maximal force scale
\newcommand{\vscore}{v_{\swirlarrow}}                    % shorthand: |v_swirl| at r=r_c
\newcommand{\vnorm}{\lVert \mathbf{v}_{\mkern-2mu\scriptscriptstyle\boldsymbol{\circlearrowleft}} \rVert}  % swirl speed magnitude
\newcommand{\rhoF}{\rho_{\!f}}\newcommand{\rhof}{\rho_{\!f}}     % effective fluid density
\newcommand{\rhoE}{\rho_{\!E}}\newcommand{\rhoe}{\rho_{\!E}}                           % swirl energy density
\newcommand{\rhoM}{\rho_{\!m}}\newcommand{\rhom}{\rho_{\!m}}                           % mass‑equivalent density
\newcommand{\omegas}{\boldsymbol{\omega}_{\swirlarrow}}  % swirl vorticity
\newcommand{\Om}{\Omega_{\swirlarrow}}                   % swirl angular frequency profile
\newcommand{\rc}{r_c}                                    % canonical horn/return-flow circulation radius
\newcommand{\rhocore}{\rho_{\text{core}}}
\newcommand{\GammaSST}{\Gamma_0}                         % legacy/core circulation quantum used in the hydrodynamic sector
\newcommand{\GammaCore}{\Gamma_{\rm core}}                    % core/horn-flow circulation
\newcommand{\GammaC}{\Gamma_C}                                % Compton streamline circulation
\newcommand{\lambdabarC}{\bar{\lambda}_C}

% additional macros for constants and notation
\newcommand{\phiG}{\varphi}
\newcommand{\me}{m_e}
\newcommand{\Eeff}{\mathcal{E}_{\rm eff}}

\titleformat{\section}{\large\bfseries}{\thesection}{1em}{}

%=========================================
% DOCUMENT
%=========================================
\begin{document}

\thispagestyle{empty}%
\begin{center}
    \Large \bfseries \papertitle \par \vspace{1cm}
    {\Large \itshape \textbf{Omar Iskandarani}\textsuperscript{\textbf{*}} \par} \vspace{0.5cm}
    {\small Independent Researcher, Groningen, The Netherlands\par} \vspace{0.5cm}
    {\today \par}  \vspace{0.5cm}
\end{center}

\begin{abstract}
    Swirl‑String Theory (SST) models elementary particles as knotted vortex loops in an incompressible, inviscid superfluid.  Unlike point‑particle or field‑theoretic descriptions, SST attributes the mass and neutral spin normalization of lepton candidates to the hydrodynamic properties of structured vorticity.  In this manuscript we derive the electron mass scale and the reduced Planck constant from a Golden‑parameterized Nonlinear Schrödinger (NLS) vortex core, while treating the visible electromagnetic magnetic moment through a two-radius gyromagnetic locking condition.  We show that the self‑intersection of a trefoil vortex ring ($T_{3,2}$) collapses into a Hill spherical vortex; this ``degenerate torus'' reproduces the observed electron rest mass.  We then compute the neutral hydrodynamic angular-momentum normalization and recover the physical value of $\hbar$ from the circulation quantum and flow-envelope density.  The Bohr magneton is not modeled as a literal electromagnetic current loop at the femtometre core scale; instead, a charged sector must lock to the Compton streamline, $R_Q v_Q=\hbar/m_e$, while the neutral hydrodynamic sector supplies $S=\hbar/2$.  Acoustic or medium dressing is retained only as a candidate bridge to the anomalous correction $a_e$, with any additional Hopfion/core form-factor contribution required to satisfy precision $g-2$ bounds.  Finally, we demonstrate that higher‑generation leptons correspond to discrete topological states of $T_{p,2}$ torus knots with masses scaling as $\varphi^{-2k}$, giving a quantitative description of the muon and tau.  Throughout, we compare our hydrodynamic quantities with textbook definitions of the classical electron radius, the Compton frequency, the Bohr magneton and the anomalous magnetic moment\cite{QEDanom}.
\end{abstract}
Keywords: \paperkeywords

\begin{textblock*}{1\textwidth}(0.15\textwidth,1.11\textheight)
    \raggedright\footnotesize
    \renewcommand{\arraystretch}{1.0}
    \noindent\rule{\textwidth}{0.4pt}\\[0.5em]
    \textsuperscript{\textbf{*}} Email: \texttt{info@omariskandarani.com}\\
    \ \ \ ORCID: \texttt{\href{https://orcid.org/0009-0006-1686-3961}{0009-0006-1686-3961}}  \\
    \ \ \ \ DOI: \href{https://doi.org/\paperdoi}{\paperdoi}
\end{textblock*}

%======================================================================
\section{Introduction and Axioms}
\label{sec:introduction}

Swirl‑String Theory (SST) is predicated on the existence of a universal, incompressible and inviscid fluid medium.  Mass and gravity arise not from curvature or quantized fields but from local vorticity and pressure gradients within this medium.  A particle is represented by a closed loop of vorticity (a swirl‑string), characterised by a finite core radius $\rc$ and circulation quantum $\GammaSST$.  To avoid the singularities inherent to thin‑tube approximations, SST adopts a Nonlinear Schrödinger (NLS) vortex ring with a ``healing length'' equal to the core radius.  The requirement that the pressure at the core boundary is constant (``const druk'') uniquely selects the Golden Ratio $\phiG=(1+\sqrt{5})/2$ as a parameter of the NLS solution.

The aim of this work is twofold: (i) derive the electron mass scale and neutral spin normalization ($m_e$, $\hbar$) from the hydrodynamics of the swirl‑string, while formulating the Bohr magneton and $g$-factor as a conditional two-sector locking result, and (ii) elucidate the discrete scale invariance that leads to the lepton mass spectrum.  Our approach is self‑contained at the level of the hydrodynamic constant chain, but the electromagnetic moment sector is explicitly status-stratified: a literal femtometre current-loop reading is rejected, while the Dirac/Compton magnetic moment is recovered if the charged sector locks to the Compton streamline.

\paragraph{Canonical constants.}  We anchor our derivations using the classical electron radius $r_e$, defined by $r_e = e^2/(4\pi\varepsilon_0 m_e c^2)$ or equivalently $r_e = \alpha\hbar/(m_e c)$, where $e$ is the elementary charge, $m_e$ the electron mass, $\varepsilon_0$ the vacuum permittivity, $c$ the speed of light and $\alpha$ the fine‑structure constant.  In the corrected reading used here, $\rc = r_e/2$ is a horn/return-flow circulation radius rather than the electromagnetic charge radius or a literal charged tube radius.  The characteristic swirl speed at this circulation radius is set by the Compton frequency $\omega_c = m_e c^2/\hbar$, leading to $\vnorm = \omega_c\,\rc$.  Equivalently, with $\lambdabarC=\hbar/(m_e c)$ and $\lambda_C=h/(m_e c)$,
\begin{equation}
    \rc = \frac{\alpha}{2}\,\lambdabarC = \frac{\alpha}{4\pi}\,\lambda_C .
    \label{eq:rc_compton_relation}
\end{equation}
We distinguish the core/horn-flow circulation from the Compton streamline circulation:
\begin{equation}
    \GammaCore = 2\pi\rc\vnorm
    = \frac{\alpha^2}{4}\,\frac{h}{m_e},
    \qquad
    \GammaC = 2\pi\lambdabarC c = \frac{h}{m_e}.
    \label{eq:gamma_core_vs_compton}
\end{equation}
The Planck constant $\hbar$ is recovered as a neutral hydrodynamic spin normalization.  The Bohr magneton $\mu_B = e\hbar/(2 m_e)$ is recovered only after imposing a separate electromagnetic locking condition on the charged sector.

\paragraph{Notation.}  We use $\rhocore$ for the core density, $\rhoF$ for the effective fluid density at large scales, and $\rhoe = \rhocore c^2$ for the energy density.  All numerical evaluations will use CODATA 2025 values unless otherwise stated.

\paragraph{Structure.}  Section~\ref{sec:NLS} formulates the Golden NLS vortex and solves for the degenerate torus representing the electron.  Section~\ref{sec:emergent_constants} derives the neutral hydrodynamic spin normalization $\hbar$ and then separates the rejected core current-loop moment from the conditional Compton-streamline route to the Bohr magneton and $g=2$.  Section~\ref{sec:leptons} relates higher‑generation leptons to topological torus knots.  Appendices provide extended derivations and discuss discrete tangle complexity.

%======================================================================
\section{Golden NLS Vortex and Degenerate Torus}
\label{sec:NLS}

The energy of a vortex ring of major radius $R$ and core radius $\rc$ in a compressible superfluid can be expressed in the NLS framework as
\begin{equation}
    E = \tfrac{1}{2} \rhocore \GammaSST^2 R \left[ \ln \left( \frac{8R}{\rc} \right) - A \right],
    \label{eq:NLS_energy}
\end{equation}
where $A$ is a dimensionless constant determined by the core pressure profile.  In SST the const‑druk condition requires that the pressure at the boundary is constant; thermodynamic analysis (see Appendix~\ref{app:golden}) shows that this is satisfied by setting $A=\phiG$ and the corresponding translational coefficient $B=\phiG^{-1}$.  The Golden Ratio therefore controls the self‑similar scaling of the vortex core.

Using mass–energy equivalence $E = m c^2$, the mass functional becomes
\begin{equation}
    m(R) = \frac{\rhocore \GammaSST^2 R}{2 c^2} \left[ \ln \left( \frac{8R}{\rc} \right) - \phiG \right].
    \label{eq:mass_functional}
\end{equation}
Given the canonical triad $(\rhocore,\GammaSST,\rc)$, we solve the transcendental equation $m(R)=m_e$ for the dimensionless ratio $x=R/\rc$.  Substituting $\rhocore = 3.893\times 10^{18}\,\text{kg m}^{-3}$, $\GammaSST = 9.6845\times 10^{-9}\,\text{m}^2\text{s}^{-1}$ and $\rc = 1.40897\times 10^{-15}\,\text{m}$, Eq.~\eqref{eq:mass_functional} yields
\begin{equation}
    x\,[\ln(8x) - \phiG] \approx 0.31826,\quad x \approx 0.895.
\end{equation}
Hence the major radius is $R_e \approx 0.895\,\rc$; it is smaller than the core radius, implying that the toroidal hole collapses and the vortex forms a contiguous spherical droplet.  This configuration corresponds exactly to Hill's spherical vortex~\cite{Hill1894}, an exact solution of the Euler equations discovered by Hill in 1894 and often used to model vortex rings.  The resulting structure is localized as a neutral hydrodynamic carrier at the scale $\rc$.  It should not be read as an electromagnetic charge distribution of radius $\rc$; pointlike electromagnetic behavior must instead be supplied by the charged Compton-streamline sector introduced below.

\subsection{Hooke closure and maximal swirl force}

The Hooke closure postulates that the maximum swirl force within the core balances a harmonic restoring force, $m_e c^2 = \tfrac{1}{2} k r_e^2$.  With $\rc = r_e/2$, the maximal swirl force is
\begin{equation}
    \Fmax = \frac{m_e c^2}{2\rc} \approx 29.05\,\text{N},
\end{equation}
and the core density is
\begin{equation}
    \rhocore = \frac{\Fmax}{\pi \rc^2 \vnorm^2},
\end{equation}
where $\vnorm = \omega_c\,\rc$ is the characteristic swirl speed at the core boundary.  These expressions are derived and numerically validated in Appendix~\ref{app:constants}.

%======================================================================
\section{Emergent Constants and Magnetic Moments}
\label{sec:emergent_constants}

Having established the degenerate torus topology, we now derive the neutral hydrodynamic angular-momentum normalization of the vortex droplet and then state the separate electromagnetic locking condition needed for the observed magnetic moment.  This separation is essential: a co-distributed charged rotor would give the classical value $g=1$, while the observed electron requires a Dirac/Compton charge sector with $g=2$ at tree level.

\subsection{Planck constant from hydrodynamic spin}

The neutral spin normalization of the droplet arises from the circulation of the core fluid.  Treating the degenerate torus as a neutral flow envelope of characteristic radius $\rc$ with density $\rhocore$, the angular-momentum scale is the product of the relevant hydrodynamic inertia and circulation.  In SST this yields
\begin{equation}
    \hbar_{\rm SST} = \rhocore \rc^3 \GammaSST.
    \label{eq:hbar_sst}
\end{equation}
Substituting numerical values gives
\begin{equation}
    \hbar_{\rm SST} = (3.893\times 10^{18})\,(2.797\times 10^{-45})\,(9.684\times 10^{-9}) \approx 1.054\times 10^{-34}\,\text{J s},
\end{equation}
in excellent agreement with the physical reduced Planck constant $\hbar$.  The physical spin assigned to the electron candidate is then
\begin{equation}
    S_H = \frac{\hbar_{\rm SST}}{2} \approx \frac{\hbar}{2}.
    \label{eq:hydro_spin_half}
\end{equation}
Thus the hydrodynamic sector supplies the spin normalization, not a literal electromagnetic current.  Quantization is interpreted as emerging from hydrodynamic stiffness and topology rather than being postulated.

\subsection{Bohr magneton and Dirac $g=2$}

A literal current loop of charge $e$ moving at the core speed on the radius $\rc$ would have
\begin{equation}
    \mu_{\rm core-loop}
    = \frac{1}{2}e\,\vnorm\rc .
    \label{eq:core_loop_mu}
\end{equation}
Using $\vnorm=\alpha c/2$ and $\rc=\alpha\lambdabarC/2$, its ratio to the Bohr magneton is
\begin{equation}
    \frac{\mu_{\rm core-loop}}{\mu_B}
    = \frac{\vnorm\rc}{c\lambdabarC}
    = \frac{\alpha^2}{4}
    \approx 1.33\times10^{-5}.
    \label{eq:core_loop_failure}
\end{equation}
This is the old literal femtometre current-loop reading.  It is rejected here as both numerically wrong and incompatible with a pointlike electromagnetic form factor.

The retained SST statement is instead a two-radius gyromagnetic locking condition.  Let the neutral hydrodynamic carrier supply the spin
\begin{equation}
    S_H=\frac{\hbar_{\rm SST}}{2},
\end{equation}
and let a charged sector move on a streamline of radius $R_Q$ with tangential speed $v_Q$.  Its magnetic moment is
\begin{equation}
    \mu_Q=\frac{1}{2}eR_Qv_Q.
    \label{eq:charge_streamline_mu}
\end{equation}
The gyromagnetic ratio defined by $\mu_Q=g(e/2m_e)S_H$ is therefore
\begin{equation}
    g_{\rm tree}
    =\frac{2m_e R_Qv_Q}{\hbar_{\rm SST}}.
    \label{eq:g_locking_general}
\end{equation}
Thus the Dirac tree-level value follows if and only if the charged sector locks to the Compton streamline,
\begin{equation}
    R_Qv_Q=\frac{\hbar}{m_e}
    \quad\Longleftrightarrow\quad
    2\pi R_Qv_Q=\GammaC=\frac{h}{m_e}.
    \label{eq:compton_streamline_lock}
\end{equation}
For the canonical Compton streamline $R_Q=\lambdabarC$ and $v_Q=c$, Eq.~\eqref{eq:compton_streamline_lock} gives
\begin{equation}
    \mu_Q=\frac{e\hbar}{2m_e}=\mu_B,
    \qquad
    g_{\rm tree}=2.
    \label{eq:bohr_from_compton_lock}
\end{equation}
Accordingly, Eq.~\eqref{eq:hbar_sst} is retained as a hydrodynamic spin normalization, while the visible electromagnetic moment is not computed from a co-rotating charged core.  The factor of two is the ratio between a full Compton circulation quantum in the charged sector and the half-quantum spin assignment in Eq.~\eqref{eq:hydro_spin_half}.  The dynamical selection of the locking condition \eqref{eq:compton_streamline_lock} is a separate variational gate, not yet a theorem.

\subsection{Acoustic dressing and the anomalous magnetic moment}

Standard QED attributes the anomalous magnetic moment $a_e=(g-2)/2$ to radiative corrections.  In the corrected SST reading, the leading magnetic moment is the Dirac/Compton result of Eq.~\eqref{eq:bohr_from_compton_lock}; any acoustic or medium dressing is an additional bridge contribution and must not generate a visible femtometre-scale form factor.  The canon-compatible target is
\begin{equation}
    a_e
    = \frac{\alpha}{2\pi} + \mathcal{O}(\alpha^2) + \delta a_H,
    \label{eq:ae}
\end{equation}
where $\delta a_H$ denotes any extra Hopfion/core contribution beyond orthodox QED-like dressing.  Precision $g-2$ agreement requires
\begin{equation}
    |\delta a_H| \ll 10^{-12}
    \label{eq:delta_aH_bound}
\end{equation}
at the level of the electron anomaly.  This is the magnetic-moment form-factor gate of the present model.

The complete magnetic moment is therefore written as
\begin{equation}
    \mu_e
    = \mu_B\left[1 + \frac{\alpha}{2\pi} + \mathcal{O}(\alpha^2) + \delta a_H\right],
    \label{eq:mu_complete_corrected}
\end{equation}
with $\mu_B=e\hbar/(2m_e)$ inherited from the Compton-streamline locking condition rather than from an electromagnetic current loop at radius $\rc$.

\subsection{Status of the gyromagnetic route}
\label{subsec:g_status}

The corrected magnetic-moment sector has the following status hierarchy.
\begin{align}
    \mu_{\rm core-loop}
    &= \frac{1}{2}e\vnorm\rc,
    &\frac{\mu_{\rm core-loop}}{\mu_B}
    &= \frac{\alpha^2}{4},
    &\text{literal core current loop} &\quad [\mathrm{rejected}],\\
    \hbar_{\rm SST}
    &= \rhocore\rc^3\GammaSST,
    &S_H &= \frac{\hbar_{\rm SST}}{2},
    &\text{neutral hydrodynamic spin} &\quad [\mathrm{retained}],\\
    R_Qv_Q
    &= \frac{\hbar}{m_e},
    &g_{\rm tree} &= 2,
    &\text{Compton-streamline locking} &\quad [\mathrm{conditional}].
\end{align}
The open variational task is to show that the bound charged sector has a stable stationary point on the streamline satisfying Eq.~\eqref{eq:compton_streamline_lock}.  If this fails by more than the precision allowed in Eq.~\eqref{eq:delta_aH_bound}, the magnetic moment sector cannot be treated as derived from the present hydrodynamic model.

%======================================================================
\section{Topological Lepton Generations}
\label{sec:leptons}

SST predicts that heavier leptons correspond to more complex torus knot topologies.  A $(p,2)$ torus knot has $p$ meridional crossings and two longitudinal windings.  The vortex energy scales linearly with $p$ and exponentially with discrete screening layers indexed by $k$; thus the mass of a lepton with knot number $p$ and layer index $k$ is
\begin{equation}
    m(p,k) = m_e \left(\frac{p}{3}\right) \phiG^{-2k}.
    \label{eq:lepton_mass}
\end{equation}
For the electron ($p=3$) we take $k=0$.  Choosing $(p,k)=(5,5)$ yields $m \approx 104.7\,\text{MeV}/c^2$, matching the muon within 1\%.  Similarly, $(p,k)=(11,7)$ gives $m \approx 1579\,\text{MeV}/c^2$, consistent with the tau to within 10\%.  Deviations reflect higher self‑interaction harmonics and are discussed in Appendix~\ref{app:knots}.

\begin{table}[h]
    \centering
    \caption{Comparison of SST lepton masses with experimental values.  Here $p$ is the torus knot crossing number and $k$ the Golden‑layer index.}
    \begin{tabular}{@{}lllll@{}}
        \toprule
        Lepton & Knot $T_{p,2}$ & $k$ & SST mass [MeV] & Physical mass [MeV] \\
        \midrule
        Electron & $T_{3,2}$ & 0 & 0.511 & 0.511 \\
        Muon     & $T_{5,2}$ & 5 & 104.7 & 105.7 \\
        Tau      & $T_{11,2}$ & 7 & 1579  & 1776.9 \\
        \bottomrule
    \end{tabular}
    \label{tab:leptons}
\end{table}

The discrete nature of $k$ reflects the allowed resonant acoustic modes within the Golden‑parameterized NLS core.  Because $k$ increases roughly linearly with the crossing number, heavy leptons appear as topological excitations of the same fundamental trefoil.

%======================================================================
\section{Conclusion}

We have shown that classical hydrodynamics, when endowed with a Golden‑ratio NLS vortex core and const‑druk constraint, yields a self‑consistent description of the neutral mass and spin-normalization sector of leptons.  The electron candidate emerges as a degenerate torus whose hydrodynamic angular-momentum scale gives rise to the reduced Planck constant.  The visible Bohr magneton is not attributed to a co-rotating charged core; it is recovered conditionally when a charged sector locks to the Compton streamline carrying $\GammaC=h/m_e$, while the neutral hydrodynamic sector supplies $S=\hbar/2$.  Discrete scale invariance explains the existence and approximate masses of the muon and tau as topological excitations.  These results suggest that quantum mechanics may be an emergent description of a knotted superfluid, provided the remaining variational locking and precision $g-2$ gates are satisfied.

\section*{Acknowledgments}

The author thanks the Swirl‑String Theory collaboration for discussions and the developers of open‑source computational tools.  Any errors or omissions remain the author's responsibility.

\appendix

\section{Golden NLS Derivation}
\label{app:golden}

The NLS vortex solution possesses two dimensionless constants $A$ and $B$ that determine the core pressure and translational velocity.  Thermodynamic stability requires that the pressure on the core boundary be uniform and finite.  Setting the net pressure gradient to zero yields $B = A - 1$.  Requiring that the constants be positive and that $B = 1/A$ (for self‑similarity) leads uniquely to $A = \phiG$ and $B = \phiG^{-1}$.  The identity $\phiG - 1 = \phiG^{-1}$ ensures that the velocity and pressure terms satisfy the const‑druk constraint.  This argument underpins the choice of $\phiG$ in Eq.~\eqref{eq:NLS_energy}.

\section{Derivation of Canonical Constants}
\label{app:constants}

Starting from the classical electron radius $r_e = \alpha\hbar/(\me c)$ and identifying $\rc = r_e/2$, the corrected interpretation is that $\rc$ is the horn/return-flow circulation radius of the neutral carrier, not the electromagnetic charge radius.  The Hooke closure assumes $m_e c^2 = \tfrac{1}{2} k r_e^2$.  The effective spring constant is therefore $k = 2 m_e c^2/r_e^2$.  The maximal swirl force is $F_{\rm swirl}^{\max} = k\,\rc = m_e c^2/(2\rc)$.  The flow-envelope density can then be expressed as
\begin{equation}
    \rhocore = \frac{F_\textrm{swirl}^{\max}}{\pi \rc^2 \vnorm^2},
    \label{eq:equation}
\end{equation}
where $\vnorm = \omega_c\,\rc$ and $\omega_c = m_e c^2/\hbar$.  Evaluating numerically gives $\rhocore \approx 3.893\times 10^{18}\,\text{kg m}^{-3}$.  In this patched reading $\rhocore$ is a flow-envelope density closure, not the density of an electromagnetic charge distribution.  The effective fluid density at large scales follows by equating the swirl energy density $\rhoe$ with the fluid kinetic energy; solving yields $\rhoF \approx 6.84\times 10^{-7}\,\text{kg m}^{-3}$, consistent with an ultra‑low‑density superfluid.

\section{Torus Knot Mass Spectrum}
\label{app:knots}

The mass spectrum of higher‑order knots can be obtained by considering the line energy per unit length (line tension) of the vortex string and summing over the total length of the knot.  For a $(p,2)$ torus knot, the minimum crossing number is $p$ and the length scales linearly with $p$.  Discrete scale invariance implies that each additional screening layer multiplies the mass by a factor $\phiG^{-2}$.  Combining these effects leads to Eq.~\eqref{eq:lepton_mass}.  A more refined model incorporating the Biot–Savart self‑interaction energy and tangle complexity is discussed in Appendix~C of the original SST manuscript; here we summarise only the leading behaviour.

%======================================================================
\begin{thebibliography}{99}

\bibitem{Hill1894} M.~J.~M. Hill, ``On a spherical vortex,'' \emph{Philosophical Transactions of the Royal Society A}, vol.~185, pp.~213--245, 1894.

\bibitem{Pethick2008} C.~J. Pethick and H.~Smith, \emph{Bose--Einstein Condensation in Dilute Gases}, Cambridge University Press, 2008.

\bibitem{Berloff2004} N.~G. Berloff, ``Interactions of vortices in a trapped Bose--Einstein condensate,'' \emph{Journal of Physics A: Mathematical and General}, vol.~37, pp.~5–19, 2004, \url{https://doi.org/10.1088/0305-4470/37/5/011}.

\bibitem{Conway1970} J.~H. Conway, ``An enumeration of knots and links, and some of their algebraic properties,'' in \emph{Computational Problems in Abstract Algebra}, Pergamon Press, 1970, pp.~329–358.

\bibitem{Kauffman2001} L.~H. Kauffman, \emph{Knots and Physics}, 3rd~ed., World Scientific, 2001.

\bibitem{QEDanom} T.~Kinoshita, ``Theory of the anomalous magnetic moment of the electron,'' \emph{Reports on Progress in Physics}, vol.~59, pp.~1459–1510, 1996.

\end{thebibliography}

\end{document}